% Cote 119: offprints, reports, courses, CV and letters (typed and printed part): archive cote 119, site page 139 (PDF page 140).
% typescript. Draft compared: none.
% AI-transcribed by Claude Opus 5.5 (claude-opus-5-5, adaptive thinking, effort high), Message Batches API, batch r119_full (msgbatch_01DwZNm5Y7YsKBWyAN4sFs8w), 2026-11-13; not yet checked by eye.
\scanpage{119/139}
Pour terminer ce rapport préliminaire,  je me sens en mesure d'appo
% NOTE: the right edge of the scan cuts off the end of this line (presumably "r-")
ter quelques précisions pratiques au programme ``tous azimuths'' exposé
d\over{\unsure{ans}}{ans} l'Esquisse d'un Programme (qui sera jointe à ``Récoltes et Semailles'',
ainsi que \inserted[typed]{le présent rapport,} l'Esquisse Thématique et quelques autres textes de nature ma-%
thématique,  pour \over{\unsure{c}}{c}onstituer le volume 1 des Réf\over{\unsure{le}}{le}xions Mathématiques).
Je tiens \inserted[typed]{d'abord,} \deleted{avant} toute autre chose,\hand{à} m'\unsure{ác|ac}quitter (comme je le dis dans
l'Esquisse d'un Programme, par.7) ``d'une dette vis-à-vis de mon passé
scientifique'',  en esquissant tout au moins \xout{\illegible{3}} \inserted[typed]{dans les grandes lignes} des
principales \xout{idxx} visions d'ensemble auxquel\unsure{les} j'étais parvenu entre 1955
et 1970 et qui, \xout{ontxété} sans avoir trouvé \unsure{alors} la forme d'exposés systé-%
matiques et publiés,  ont été laissés pour compte par mes ex-élèves
après mon ``départ'' en 1970.  Il s'agit d'une part de l'ensemble d'intui-%
tions et d'idées en direction de l'``algèbre topologique''\hand{,} à laquelle j'ai
fait allusion tantôt,  et d'autre part  d'un ``vaste tableau des motifs''
qu'il s'agit de\hand{``}brosser à grands traits\hand{''},  apte à servir à la fois de
maître d'oeuvre pour l'édification d'une théorie qui reste conjecturale,
et de fil conducteur très sûr et \deleted{d'outil d'}\inserted{de fécond} instrument de découverte,
pour pouvoir prédire ``à quoi on est en droit de s'attendre'' dans une
foule de situations impliquant les propriétés géométrico-arithmétiques
de la cohomologie des variétés algébriques. (Une version très parcellai
% NOTE: the right edge of the scan cuts off the end of this line (presumably "re")
de certaines de mes principales idées à de sujet est présentée,  sans
que mon nom y soit prononcé,  dans le volume ``Hodge Cycles, Motives, and
Shimura Varieties'' des Lecture Notes (n\textsuperscript{o} 900), \inserted[typed]{1982,} par Pierre Deligne, James
S. Milne, Arthur Ogus, Kuang-yen Shih.) Je compte poursuivre cette parti
% NOTE: the right edge of the scan cuts off the end of this line (presumably "e")
de mon programme, au cours des deux ou trois années qui suivent,  dans
les deux ou trois volumes des Réflexions Mathématiques faisant suite aux
deux volumes dont je suis en train de terminer la préparation.  Ceci
fait,  je prévois de me consacrer prioritairement au programme de géomé-%
trie algébrique anabélienne -  d'une part,  tracer à grands traits les
principales idées,\xout{et} conjectures et résultats déjà obtenus,  et d'autre
part,  entreprendre une étude gémétrico-arithmétique minutieuse de la
``tour de Teichmüller'',  et plus particulièrement,  de ses deux premiers
étages.

\xout{\illegible{1}} le 10.12.1984

\hand{Alexandre Grothendieck}

\mspage{139}
